\chapter{Introduction}\label{sec:intro}

The Sherrington--Kirkpatrick (SK) model is the mean-field Ising spin glass,
in which every pair of spins interacts through an independent Gaussian
coupling~\cite{sk-original}. Its equilibrium theory is well developed.
Parisi's hierarchical ansatz~\cite{parisi-1979} predicted the limiting free
energy, and this prediction was proved through Guerra's interpolation
bound~\cite{guerra-2003} and Talagrand's matching lower
bound~\cite{parisi-formula}. At every fixed inverse temperature $\beta>1$
and zero external field, the minimizing Parisi measure is supported on an
interval $[0,q_\beta]$; this was proved by Zhou for $\beta$ close to
one~\cite[Theorem~1]{zhou-2026} and in~\cite[Theorem~1.1]{full-rsb} for every
$\beta>1$. Much less is known about the
dynamics of the model at low temperature, and about the finite-size
corrections to its equilibrium quantities. This paper addresses both.

Our main results concern the heat-bath dynamics, also called Glauber
dynamics~\cite{glauber-1963}, in which one spin at a time is resampled from
its conditional law. For fixed $\beta>1$, we sample the couplings and then one
configuration from the Gibbs measure, and run the dynamics from this
configuration. We prove that, with probability tending to one, the law of
the configuration at time $\exp(A_\beta n^{7/30})$ is still at total-variation
distance greater than $\tau$ from the Gibbs measure, for every fixed $\tau<1/2$
and a coefficient $A_\beta>0$ that does not depend on $\tau$
(\cref{thm:barrier-730}). In the other direction, for every fixed
$\omega\in(0,1/4)$, the distance is at most $1/4$ at time $\exp(n^{1-\omega})$,
again with probability tending to one (\cref{thm:upper}). OpenAI proved bounds
of both kinds in the same setting, with the same joint law of the couplings
and the initial configuration, the same clocks of rate one per site, and the
threshold $1/4$: a lower bound at time
$\exp(n^{1/10000})$~\cite[Theorem~1.1]{oai-barriers} and an upper bound at
time $\exp(n^{1-1/40000000})$~\cite[Theorem~1]{oai-typical-upper}. We follow
the overall strategy of both proofs. Let $T^J(X_0)$ be the first time at which
the distance from the Gibbs measure is at most $1/4$. Combining our two
bounds, we find that if $\log T^J(X_0)=n^{\zeta+o(1)}$, then
$7/30\le\zeta\le3/4$ (\cref{cor:two-sided}). The proof of the lower bound
uses the upper bound $n^{-7/10}$ of Aronow and Lopatto on the mean bias of
the pressure~\cite[Theorem~1.1]{hexagon-bias}.

Our second contribution is a complete formal proof of
\cref{thm:barrier-730} in the Lean~4 proof assistant~\cite{lean4,mathlib},
together with formal proofs of two results that are not new. The first is the
fluctuation theorem $\Var(F_n)=n^{1/3+o(1)}$ for the logarithm $F_n$ of the
partition function (\cref{f:main}), which was stated and proved by OpenAI
in~\cite[Theorem~1.1]{oai-fluctuations}. OpenAI's public formalization of
this theorem assumes a two-sided bound on positive logarithmic moments of
the centered completed log partition function, which is proved on paper
in~\cite[Proposition~7.5]{oai-fluctuations}; our formal proof establishes
this bound. We restate the theorem with attribution,
describe the structure of its proof, and reproduce the short final step of
that proof; we do not claim a new mathematical proof of \cref{f:main}.
\statusIntroShort

The second known result is a lower bound with exponent $5/6$ on the mean bias
$P_\beta-n^{-1}\E F_n$, where $P_\beta=\lim_{n\to\infty}n^{-1}\E F_n$ is the
limiting pressure. It follows from \cref{f:main} and a variance--bias
inequality of Aronow and Lopatto
\cite[Lemma~7.4 and Remark~7.14]{quantitative-parisi} that the mean bias is
at least of order $n^{-5/6-\varepsilon}$ for every fixed $\varepsilon>0$
(\cref{cor:mean-bias}). This deduction was already recorded
in~\cite[Section~1.6]{hexagon-bias}; what is new here is its formal proof.
Together with the upper bound of~\cite{hexagon-bias}, it shows that if the
mean bias is $n^{-a+o(1)}$, then $7/10\le a\le5/6$.

\section{Main results}\label{sec:intro-results}

For an integer $n\ge2$, let $\Sigma_n=\{-1,1\}^n$, and let
$J=(J_{ij})_{1\le i<j\le n}$ be independent standard Gaussian variables. We
set $J_{ji}=J_{ij}$ and $J_{ii}=0$. The Hamiltonian, the partition function,
and the Gibbs measure at inverse temperature $\beta$ are
\begin{equation}\label{eq:intro-model}
 H_n(x)=\frac1{\sqrt n}\sum_{1\le i<j\le n}J_{ij}x_ix_j,\qquad
 Z_n=\sum_{x\in\Sigma_n}e^{\beta H_n(x)},\qquad
 \pi^J(x)=\frac{e^{\beta H_n(x)}}{Z_n}.
\end{equation}
We write $R(x,y)=n^{-1}\sum_{i=1}^n x_iy_i$ for the overlap of
$x,y\in\Sigma_n$, and $F_n=\log Z_n$.

In the heat-bath dynamics, each site carries an independent exponential
clock of rate one. When the clock of site $i$ rings at a configuration $x$,
the spin $x_i$ is replaced by $\sigma\in\{-1,1\}$ with probability
\[
 p_i^J(\sigma\mid x)=\frac{e^{\beta\sigma h_i(x)}}{2\cosh(\beta h_i(x))},\qquad
 h_i(x)=\frac1{\sqrt n}\sum_{j\ne i}J_{ij}x_j,
\]
which is the conditional law of the $i$th spin under $\pi^J$ given the
other spins. An update may leave the configuration unchanged. We write
$P_t^J(x,\cdot)$ for the law at time $t$ of this process started at $x$, and
$K^J$ for the transition matrix of the discrete chain that chooses a site
uniformly at random and performs one such update. One unit of continuous
time contains $n$ update attempts on average. Both chains are reversible with
respect to $\pi^J$. For $x\in\Sigma_n$, $t\ge0$, and integers $k\ge0$, we
set
\[
 d^J(x,t)=\norm{P_t^J(x,\cdot)-\pi^J}_{\TV},\qquad
 d^J_{\mathrm{disc}}(x,k)=\norm{(K^J)^k(x,\cdot)-\pi^J}_{\TV},
\]
where $\norm{\nu-\zeta}_{\TV}=\frac12\sum_x\abs{\nu(x)-\zeta(x)}$. We also
write
\[
 T^J(x)=\inf\{t\ge0:d^J(x,t)\le1/4\},\qquad
 T^J_{\mathrm{disc}}(x)=\min\{k\ge0:d^J_{\mathrm{disc}}(x,k)\le1/4\}
\]
for the mixing times of the two chains started at $x$. Since
$t\mapsto d^J(x,t)$ is continuous and nonincreasing
(\cref{lem:model-dynamics}), we have $T^J(x)\le t$ if and only if
$d^J(x,t)\le1/4$, and similarly in discrete time. Finally,
$\PP_{J,\,X_0\sim\pi^J}$ denotes the joint law of the couplings $J$ and of a
configuration $X_0$ whose conditional law given $J$ is $\pi^J$; see
\eqref{eq:common-joint}. These definitions are given in detail in
\cref{sec:model}.

\begin{theorem}[Barrier from typical equilibrium states]\label{thm:barrier-730}
For every fixed $\beta>1$ there exists $A_\beta>0$ such that, for every
fixed $\tau\in(0,1/2)$,
\begin{align}
 \PP_{J,\,X_0\sim\pi^J}
   \bigl\{d^J\bigl(X_0,e^{A_\beta n^{7/30}}\bigr)>\tau\bigr\}
   &\longrightarrow1,\label{eq:barrier-730-continuous}\\
 \PP_{J,\,X_0\sim\pi^J}
   \bigl\{d^J_{\mathrm{disc}}\bigl(X_0,\lfloor e^{A_\beta n^{7/30}}\rfloor\bigr)>\tau\bigr\}
   &\longrightarrow1.\label{eq:barrier-730-discrete}
\end{align}
\end{theorem}

The configuration $X_0$ is sampled once and then held fixed, so that
$d^J(X_0,t)$ measures the distance from equilibrium of the chain started at
this particular configuration. If the initial configuration were instead
averaged over $\pi^J$ before the distance is taken, the law at time $t$
would be $\pi^J$ by stationarity, and the distance would vanish. Since the
Gibbs mass of the set of initial configurations in the event of
\eqref{eq:barrier-730-continuous} takes values in $[0,1]$ and has expectation
equal to the displayed probability, \eqref{eq:barrier-730-continuous} is
equivalent to the convergence of this Gibbs mass to one in probability over
the disorder; this observation is made after
\cite[Theorem~1.1]{oai-barriers} and is recalled in \cref{sec:model}. The
coefficient $A_\beta$ depends only on $\beta$: the same coefficient serves
for every fixed threshold $\tau<1/2$ and for both clocks, while the rate of
convergence in \eqref{eq:barrier-730-continuous} and
\eqref{eq:barrier-730-discrete} may depend on $\tau$. The restriction
$\tau<1/2$ comes from spin-flip symmetry: for every configuration $v$, the
half-space $\{y:R(v,y)\le0\}$ has Gibbs mass at least $1/2$, and the proof
shows that, for a suitable configuration $v$ with $R(v,X_0)$ bounded away
from zero, the chain started at $X_0$ is unlikely to enter this half-space
before the stated time.

The exponent $7/30$ is one third of the exponent $7/10$ in the bound
$0\le P_\beta-n^{-1}\E F_n\le C_b n^{-7/10}$ on the mean bias
proved in~\cite[Theorem~1.1]{hexagon-bias}, where
we recall that $P_\beta$ is the limiting pressure. The proof needs overlap constraints on a scale
$\rho$ whose pressure deficit, of order $\rho^3$ per spin, dominates the
mean bias, and our counting argument allows a time $t$ with $\log t$ of
order $1/\rho$. With the bound $n^{-7/10}$, we can take
$\rho\asymp n^{-7/30}$. \Cref{sec:intro-methods} explains this balance in
more detail. A union bound over polynomially many constraints in the
argument requires $n\rho^4$ to be large compared with $\log n$, so the
present argument gives no exponent of $1/4$ or more, whatever the upper
bound on the mean bias. \Cref{thm:barrier-730} is proved in
\cref{sec:barrier-proof}.

Since total variation contracts under a Markov kernel that preserves
$\pi^J$, the distances $d^J(x,t)$ and $d^J_{\mathrm{disc}}(x,k)$ are
nonincreasing in $t$ and $k$. For every fixed $A>0$ and $\kappa<7/30$, we
have $n^\kappa<An^{7/30}$ for all sufficiently large $n$.
\Cref{thm:barrier-730} implies the following statement with
coefficient one, which is also proved in \cref{sec:barrier-proof}.

\begin{corollary}[Unit coefficient below the exponent $7/30$]\label{b:cor:subcritical}
For every fixed $\beta>1$, every fixed $\kappa\in(0,7/30)$, and every
fixed $\tau\in(0,1/2)$,
\[
 \PP_{J,\,X_0\sim\pi^J}\bigl\{d^J\bigl(X_0,e^{n^\kappa}\bigr)>\tau\bigr\}
   \longrightarrow1,\qquad
 \PP_{J,\,X_0\sim\pi^J}
   \bigl\{d^J_{\mathrm{disc}}\bigl(X_0,\lfloor e^{n^\kappa}\rfloor\bigr)>\tau\bigr\}
   \longrightarrow1.
\]
\end{corollary}

The corollary does not assert coefficient one at the exponent $7/30$
itself.

The proof gives more than the convergence of the probabilities in
\cref{thm:barrier-730}; we record the stronger forms in
\cref{thm:barrier-strong}. There exist $A>0$ and, for each $n$, a disorder
event of probability at least $1-Ce^{-cn^{1/15}}$ on which the following
holds for all sufficiently large $n$. The set of initial configurations $x$
with $d^J(x,t)\ge1/2-Ce^{-cn^{7/30}}$ for every $t\le e^{An^{7/30}}$, and
with the analogous bound for the discrete chain, has Gibbs mass at least
$1-Ce^{-cn^{7/30}}$. The same holds for the Metropolis dynamics and, more
generally, for every single-site dynamics with uniformly bounded rates that
leaves $\pi^J$ stationary. For every reversible such dynamics, the
relaxation time is at least $e^{An^{7/30}}/(Cn^{7/30})$. Finally, the Gibbs
mass of the slow initial configurations tends to one almost surely along the
full sequence of system sizes, however the disorders of different sizes are
coupled. The main change in the proof is a lower bound on $F_n$ from the
small-deviation inequality of Paouris and Valettas for convex functions of
Gaussian vectors~\cite[Theorem~2.1 and
Remark~2.4(3)]{paouris-valettas-2018}, in place of Chebyshev's inequality.
We have formalized \cref{thm:barrier-730} and \cref{b:cor:subcritical} in
Lean, but not \cref{thm:barrier-strong}; see \cref{sec:intro-formal}.

\begin{theorem}[Mixing from typical equilibrium states]\label{thm:upper}
For every fixed $\beta>1$ and every fixed $\omega\in(0,1/4)$,
\[
 \PP_{J,\,X_0\sim\pi^J}\bigl\{d^J\bigl(X_0,e^{n^{1-\omega}}\bigr)\le 1/4\bigr\}\longrightarrow1;
 \quad\text{equivalently}\quad
 \PP_{J,\,X_0\sim\pi^J}\bigl\{T^J(X_0)\le e^{n^{1-\omega}}\bigr\}\longrightarrow1 .
\]
\end{theorem}

As in \cref{thm:barrier-730}, the configuration $X_0$ is sampled once and
then held fixed when the distance $d^J(X_0,t)$ is formed, and the clocks
have rate one per site, as in \cref{sec:model-dynamics}. By the argument of
\cref{lem:model-joint}, applied to the complementary events
$\{d^J(X_0,t)>1/4\}$, \cref{thm:upper} is equivalent to the convergence to
zero, in probability over the disorder, of the Gibbs mass of the initial
configurations $x$ with $d^J(x,e^{n^{1-\omega}})>1/4$. The theorem asserts
nothing at $\omega=1/4$. It improves the exponent $1-1/40000000$
of~\cite[Theorem~1]{oai-typical-upper}, and it is proved in
\cref{part:upper}. \statusIntroShort

For the discrete chain, we prove the same bound after
$\lfloor e^{n^{1-\omega}}\rfloor$ update attempts. Since one unit of
continuous time contains $n$ update attempts on average, this number of
attempts corresponds to the time $e^{n^{1-\omega}}/n$. Together with
\cref{thm:barrier-730}, we obtain two-sided bounds in both time
conventions.

\begin{corollary}[Discrete time]\label{cor:upper-discrete}
For every fixed $\beta>1$ and $\omega\in(0,1/4)$,
\[
\PP_{J,X_0\sim\pi^J}\{d^J_{\mathrm{disc}}(X_0,\lfloor e^{n^{1-\omega}}\rfloor)\le1/4\}\to1.
\]
\end{corollary}

\begin{corollary}[Two-sided typical-start bounds]\label{cor:two-sided}
For every fixed $\beta>1$ and $\omega\in(0,1/4)$, with $A_\beta$ from \cref{thm:barrier-730},
\begin{gather*}
\PP_{J,X_0\sim\pi^J}\bigl\{e^{A_\beta n^{7/30}}<T^J(X_0)\le e^{n^{1-\omega}}\bigr\}\to1,\\
\PP_{J,X_0\sim\pi^J}\bigl\{\lfloor e^{A_\beta n^{7/30}}\rfloor<T^J_{\mathrm{disc}}(X_0)\le \lfloor e^{n^{1-\omega}}\rfloor\bigr\}\to1 .
\end{gather*}
\end{corollary}

In the notation of \cref{sec:model-asymptotics}, \cref{cor:two-sided}
implies that if $\log T^J(X_0)=n^{\zeta+o_\PP(1)}$ under
$\PP_{J,\,X_0\sim\pi^J}$ for some $\zeta$, then $7/30\le\zeta\le3/4$, and the
same holds for $T^J_{\mathrm{disc}}(X_0)$. The value $\zeta=1/3$ suggested by
the physics literature (\cref{sec:intro-dynamics}) lies in this interval.
\Cref{cor:upper-discrete,cor:two-sided} are proved in \cref{u:sec:proof}.

Both arguments stop at the exponent $1/4$. The proof of
\cref{thm:barrier-730} needs $n\rho^4/\log n\to\infty$ for the overlap scale
$\rho$, and the proof of \cref{thm:upper} needs the chain of inequalities
$\omega<\mathsf q<\mathsf s_{\rm b}<\mathsf s_{\rm seq}<\mathsf e<\mathsf u_{\rm E}<\mathsf k_{\rm f}/2\le1/4$
of \cref{u:sec:budget}. We do not know whether this coincidence has a common
explanation.

\section{Two known results: fluctuations and the mean bias}\label{sec:intro-known}

The other two results that we present concern the fluctuations of $F_n$ and
the mean bias. The first was proved by OpenAI. The
paper~\cite{oai-fluctuations} uses the uniform probability measure on
$\Sigma_n$ instead of the counting measure in \eqref{eq:intro-model}. This
changes $F_n$ by the constant $n\log2$ and leaves the following statement
unchanged (see \cref{sec:model-dictionary}).

\begin{theorem}[{\cite[Theorem~1.1]{oai-fluctuations}}]\label{f:main}
For every fixed $\beta>1$ and every $\eta>0$, there exists $n_0(\beta,\eta)$
such that, for all $n\ge n_0$,
\begin{align}
 n^{1/6-\eta}&\le\sqrt{\Var(F_n)}\le n^{1/6+\eta},\label{f:sd}\\
 n^{1/3-\eta}&\le\Var(F_n)\le n^{1/3+\eta}.\label{f:variance}
\end{align}
Moreover, for every $\eta>0$,
\begin{equation}\label{f:typical}
 \PP\bigl\{n^{1/6-\eta}\le\abs{F_n-\E F_n}\le n^{1/6+\eta}\bigr\}
 \longrightarrow1.
\end{equation}
\end{theorem}

The proof of \cref{f:main} in~\cite{oai-fluctuations} reduces it to a
two-sided bound on positive logarithmic moments of the centered completed
log partition function, which we state as \cref{f:window} and discuss in
\cref{sec:intro-fluct}. Our formal proof of
\cref{f:main} includes a formal proof of this bound
(\cref{sec:intro-formal}).

Let $p_n(\beta)=n^{-1}\E F_n$, so that the mean bias is
$\Delta_{n,\beta}=P_\beta-p_n(\beta)$. By Guerra's
bound~\cite[Theorem~3]{guerra-2003} and the Parisi
formula~\cite[Theorem~1.1]{parisi-formula}, we have $\Delta_{n,\beta}\ge0$
for every $n$. A lower bound on the variance implies a lower bound on the
mean bias. Aronow and Lopatto proved that for every fixed $\beta>1$ there
exists $c_\beta>0$ such that
\begin{equation}\label{eq:intro-variance-deficit}
 \Delta_{n,\beta}\ge c_\beta\left(\frac{\Var(F_n)+\beta^2/2}{n}\right)^{5/4}
\end{equation}
for every $n\ge2$~\cite[Lemma~7.4 and Remark~7.14, equation~(7.68)]{quantitative-parisi};
we state it as \cref{prop:bias-variance-deficit}. The term $\beta^2/2$ is
the variance of the Gaussian variable that completes the covariance of
$\beta H_n$ (\cref{sec:model-completion}). Inserting \eqref{f:variance}
into \eqref{eq:intro-variance-deficit}, we obtain the following corollary,
which is already recorded in~\cite[Section~1.6, p.~11]{hexagon-bias}.

\begin{corollary}[Mean bias]\label{cor:mean-bias}
For every fixed $\beta>1$ and every $\varepsilon>0$, there exist
$c_{\beta,\varepsilon}>0$ and $n_0(\beta,\varepsilon)$ such that
\begin{equation}\label{eq:mean-bias-lower}
 \Delta_{n,\beta}\ge c_{\beta,\varepsilon}n^{-5/6-\varepsilon}
 \qquad(n\ge n_0(\beta,\varepsilon)).
\end{equation}
Equivalently, $nP_\beta-\E F_n\ge c_{\beta,\varepsilon}n^{1/6-\varepsilon}$ for
$n\ge n_0(\beta,\varepsilon)$.
\end{corollary}

Together with the upper bound of~\cite{hexagon-bias}, \cref{cor:mean-bias}
implies that if $\Delta_{n,\beta}=n^{-a+o(1)}$ for some $a$, then
$7/10\le a\le5/6$. \Cref{sec:bias} proves \cref{cor:mean-bias}, which does
not use \cref{thm:barrier-730}, and discusses what it leaves open. If
$\liminf_{n\to\infty}n^{-1/3}\Var(F_n)>0$, the same inequality implies that
$\Delta_{n,\beta}\ge c_\beta n^{-5/6}$ for all sufficiently large $n$, as
observed in~\cite[Section~1.6, equation~(1.4)]{hexagon-bias}; this
hypothesis follows from the convergence of $n^{-1/3}\Var(F_n)$ to a positive
constant, which is stated in~\cite[Theorem~1.2]{oai-limiting-law}. Our formal
proof of \cref{cor:mean-bias} combines the formal proof of \cref{f:main}
with a formal proof of \eqref{eq:intro-variance-deficit}.

\section{Related results on the dynamics}\label{sec:intro-dynamics}

\paragraph{Physics.}
The slow relaxation of finite SK systems has been studied in physics since
the early 1980s. Sompolinsky and Zippelius developed the dynamical mean-field
theory of the spin-glass
phase~\cite{sompolinsky-zippelius-1981,sompolinsky-zippelius-1982}, in which
the relaxation times of a finite system diverge with its size $n$. In
simulations of the Ising SK model, Mackenzie and Young found two families of
time scales $\tau_n$. For observables invariant under global spin reversal,
$\log\tau_n$ grew like $n^{1/4}$, and for the reversal itself like
$n^{1/2}$~\cite{mackenzie-young-1982,mackenzie-young-1983}. Later work
revised this picture. Vertechi and Virasoro computed energy barriers between
metastable states numerically and found growth like $n^\alpha$ with
$\alpha=0.34\pm0.08$~\cite{vertechi-virasoro-1989}. Colborne's Monte Carlo
study gave the exponent $0.25\pm0.05$ for the barrier that governs the
ergodic time~\cite{colborne-1990}. Kinzelbach and Horner predicted
free-energy barriers of order $n^{1/3}$ from a resummed perturbative
analysis of the dynamics of finite systems~\cite{kinzelbach-horner-1991}.
The simulations of Billoire and Marinari~\cite{billoire-marinari-2001}, of
Bittner and Janke~\cite{bittner-janke-2006}, and of
Billoire~\cite{billoire-2010} are consistent with $\log\tau_n\propto n^{1/3}$,
and so is the computation of spectral gaps of small systems by Monthus and
Garel~\cite{monthus-garel-2009}. Aspelmeier and Moore studied barriers
between solutions of the Thouless--Anderson--Palmer
equations~\cite{thouless-anderson-palmer-1977,bray-moore-1981}. They found
that barriers of order $n^{1/3}$ are one of several possibilities below a
threshold for the free-energy density~\cite{aspelmeier-moore-2022}.

These works concern different quantities: energy barriers, free-energy
barriers along an overlap coordinate, autocorrelation times, and spectral
gaps. They use binary as well as Gaussian couplings, Metropolis as well as
heat-bath updates, and moderate system sizes. They suggest that the logarithm
of the relaxation time from a typical configuration grows like $n^{1/3}$,
but they do not determine a mixing exponent. A different question concerns
times that stay bounded as $n$ grows. When $n\to\infty$ is taken first, the
dynamics from a random configuration ages and does not reach
equilibrium~\cite{cugliandolo-kurchan-1994}; we refer
to~\cite[Section~3.7]{bouchaud-cugliandolo-kurchan-mezard-1997} for a
discussion of the two orders of limits.

\paragraph{High and critical temperature.}
Rigorous results on the dynamics depend strongly on the temperature. At high
temperature, functional inequalities imply fast mixing from every initial
configuration. Bauerschmidt and Bodineau proved a logarithmic Sobolev
inequality, uniform in $n$, for a Dirichlet form of gradient type under a
condition on the spectral width of the interaction matrix, which applies to
SK when $\beta<1/4$~\cite[Theorem~1 and Corollary~2]{bauerschmidt-bodineau-2019}.
Eldan, Koehler, and Zeitouni proved a Poincar\'e inequality for the
heat-bath Dirichlet form itself under a spectral condition of the same kind,
and deduced a polynomial bound on the mixing time from every initial
configuration~\cite[Theorems~1 and~11]{eldan-koehler-zeitouni-2022}. Anari,
Jain, Koehler, Pham, and Vuong obtained mixing in $O(n\log n)$ updates in the
same range through entropic
independence~\cite[Theorem~12]{anari-jain-koehler-pham-vuong-2022}. Wang
extended mixing in $O(n\log n)$ updates to every $\beta<1/2$
\cite[Theorem~1.1]{wang-2026}, and Boban, Li, and Oveis Gharan proved
polynomial mixing for $\beta<1/2+\varepsilon_0$, for an absolute constant
$\varepsilon_0>0$~\cite[Theorem~1.5 and Corollary~1.6]{boban-li-oveis-gharan-2026}.
With external fields sufficiently large at every site and their average
bounded by a fixed constant $H$, Paliy and Sohn proved a modified
logarithmic Sobolev inequality and mixing in $O_{\beta,H}(n\log n)$ update
attempts at every fixed $\beta$, with probability tending to
one~\cite[Corollary~1.5]{paliy-sohn-2026}.
For every fixed $\beta<1$, OpenAI proved that the spectral gap of the
continuous-time chain is bounded below by a positive constant with
probability tending to one~\cite[Theorem~1.1]{oai-spectral-gap}, and that
the chain started from the worst initial configuration exhibits cutoff at
time $(\log n)/(2\lambda(\beta))$, for a deterministic
$\lambda(\beta)>0$~\cite[Theorem~1.1]{oai-cutoff}. At the critical
temperature $\beta=1$, OpenAI proved that, for every deterministic $t_n$
with $t_n=o(n^{2/3})$, the Gibbs mass of initial configurations $x$ with
$d^J(x,t_n)>1/4$ tends to one in probability~\cite[Theorem~1.1]{oai-critical-slowing},
and that the mixing time is $n^{2/3+o(1)}$~\cite[Theorem~1.1]{oai-critical-mixing}.

\paragraph{Low temperature.}
At low temperature, Ben Arous and Jagannath showed that a free-energy barrier
in the overlap of two replicas forces an exponentially small spectral gap
in mean-field spin glasses~\cite[Theorem~1.3]{ben-arous-jagannath-2018}.
Their sufficient conditions for such a barrier include the vanishing of the
second derivative at zero of the covariance
function~\cite[Theorem~1.9]{ben-arous-jagannath-2018}, so they do not apply
to the quadratic covariance of SK. For the pure spherical $p$-spin models
with $p\ge3$, Gheissari and Jagannath proved that the spectral gap of the
Langevin dynamics is exponentially small at sufficiently low
temperature~\cite[Theorem~1]{gheissari-jagannath-2019}. For SK, Sellke proved
that there exist absolute constants $\beta_*,c>0$ such that, for every
$\beta>\beta_*$, the mixing time from the worst initial configuration is at
least $e^{c\beta n}$ with probability tending to
one~\cite[Theorem~1.1]{sellke-2025}. The slow initial configurations in his
proof are configurations at which almost every spin is aligned with its
local field by a uniform margin. He notes after his theorem that the chain may mix much
faster from a random initial configuration. El Alaoui, Montanari, and Sellke
proved that for $\beta>1$ no family of sampling algorithms that is stable
under small perturbations of the disorder approximates the Gibbs measure in
normalized Wasserstein distance~\cite[Theorem~2.6]{el-alaoui-montanari-sellke-sampling}.
Neither result determines the Gibbs mass of the initial configurations from
which the heat-bath chain mixes slowly.

OpenAI proved a lower bound and an upper bound for initial configurations
sampled from the Gibbs measure that hold throughout $\beta>1$.
By~\cite[Theorem~1.1]{oai-barriers}, the conclusions of \cref{thm:barrier-730}
hold with the time $\exp(n^{1/10000})$ and the threshold $\tau=1/4$. In a
companion preprint, OpenAI proved that
$\PP_{J,\,X_0\sim\pi^J}\{T^J(X_0)\le\exp(n^{1-c_\beta})\}\to1$ for some
$c_\beta\in(0,1)$, and the proof permits
$c_\beta=1/40000000$~\cite[Theorem~1]{oai-typical-upper}. Combining the two
bounds, and using that $t\mapsto d^J(x,t)$ is continuous and nonincreasing,
they obtained a two-sided bound on $T^J(X_0)$ with these
exponents~\cite[Corollary~2]{oai-typical-upper}. The upper bound
of~\cite{oai-typical-upper} is stated only in continuous time.

\Cref{thm:barrier-730} improves the exponent $1/10000$ of~\cite{oai-barriers}
to $7/30$, and it holds for every fixed threshold below $1/2$ with a
coefficient that does not depend on the threshold. \Cref{thm:upper} improves
the exponent $1-1/40000000$ of~\cite{oai-typical-upper} to $1-\omega$ for
every fixed $\omega<1/4$, and \cref{cor:upper-discrete} gives the same bound
for the discrete chain. For $\beta>\beta_*$, we find from \cref{thm:upper}
that the mixing time from a typical initial configuration is much smaller
than the lower bound
$e^{c\beta n}$ of~\cite{sellke-2025} on the mixing time from the worst initial
configuration. Our results do not determine the exponent of the typical
mixing time, which by \cref{cor:two-sided} lies in $[7/30,3/4]$ if it exists.
They are compatible with the value $1/3$ suggested by the physics literature,
but they do not decide whether this value is correct. The value $1/3$ also
appears as the exponent of $\Var(F_n)$ in \cref{f:main}; we know of no
argument relating the two exponents. The coefficient $A_\beta$ depends on
$\beta$, and our proof does not control it as $\beta\downarrow1$ or
$\beta\to\infty$. Similarly, the rate of convergence in \cref{thm:upper}
depends on $\beta$ and $\omega$, and we assert no uniformity in either
parameter.

\section{The fluctuation scale and the mean bias}\label{sec:intro-fluct}

\paragraph{Physics.}
The prediction of the fluctuation scale of $F_n$ has a separate history in
physics. In a truncated replica theory near the critical temperature, Kondor
found that the replica-normalized free energy depends on a small positive
replica number $s$ first at order $s^5$~\cite{kondor-1983}. Crisanti,
Paladin, Sommers, and Vulpiani interpreted the replica number as a tilt of
the disorder distribution, and predicted that the free energy per spin
fluctuates on the scale $n^{-5/6}$ in the low-temperature phase at zero
field~\cite{crisanti-paladin-sommers-vulpiani-1992}. At fixed $\beta$, this
is the scale $n^{1/6}$ for $F_n$. Parisi and Rizzo extended the fifth-order
calculation throughout the low-temperature phase and computed the resulting
large deviations of the free energy~\cite{parisi-rizzo-2008,parisi-rizzo-2009}.
Passing from these large deviations to the width of typical fluctuations
requires matching the tail of the distribution to its center. This is an
additional assumption, which they discuss
in~\cite[Section~VI]{parisi-rizzo-2009}; the limits $n\to\infty$ and
$s\downarrow0$ cannot simply be interchanged.
Aspelmeier, Billoire, Marinari, and Moore argued that a system of $n$ spins
resolves about $n^{1/6}$ levels of replica symmetry breaking. From this they
derived both a variance of order $n^{1/3}$ for the total free energy and a
correction of order $n^{-2/3}$ to its mean per
spin~\cite[Section~2]{aspelmeier-billoire-marinari-moore-2008}. Aspelmeier
derived exact identities relating the variance of the free energy to
overlaps between systems with correlated disorder~\cite{aspelmeier-2008}. At
zero temperature, numerical studies of the width of the distribution of the
ground-state energy have produced exponents between $3/4$ and $5/6$ for the
standard deviation per spin, and the fits are sensitive to corrections to
scaling~\cite{bouchaud-krzakala-martin-2003,boettcher-2005,boettcher-2010,palassini-2008}.

\paragraph{Rigorous results.}
Rigorous results on the fluctuations of $F_n$ depend on the temperature and
the field. For $0<\beta<1$ at zero field, $F_n$ has Gaussian fluctuations of
order one, as shown by Aizenman, Lebowitz, and Ruelle~\cite{alr-1987}. Comets
and Neveu studied these fluctuations by stochastic
calculus~\cite{comets-neveu-1995}. At $\beta=1$ the variance grows
logarithmically. Chen and Lam proved an upper bound of order
$(\log n)^2$~\cite[Theorem~1.1]{chen-lam-2019}, Schertzer proved the upper
bound $(\log n)/4+C$ and a lower bound of order
$\log\log\log n$~\cite[Theorems~1 and~2]{schertzer-2026}, Tierz bounded the
variance for every $n$ in terms of the spin-glass
susceptibility~\cite[Theorem~1.2]{tierz-critical-2026}, and Du and Huang
proved that $\Var(F_n)=(\log n)/6+O(1)$ and that the fluctuations have a
Gaussian limit~\cite[Theorem~1.3]{du-huang-2026}. Dey and Kang obtained the
variance asymptotics and a Gaussian limit along sequences $\beta_n\uparrow1$
for which $n^{1/3}(1-\beta_n^2)$ converges to a limit in
$(0,\infty]$~\cite[Theorem~1.1]{dey-kang-2026}. With a nonzero external field,
Chen, Dey, and Panchenko proved that the variance is of order $n$ and that a
central limit theorem
holds~\cite[Theorems~1 and~2]{chen-dey-panchenko-2017}. The spin space also
matters: for the spherical SK model at low temperature, Baik and Lee proved
that the free energy per spin has Tracy--Widom fluctuations of order
$n^{-2/3}$~\cite[Theorems~2.10 and~2.11]{baik-lee-2016}, so that the total
free energy fluctuates on the scale $n^{1/3}$ rather than $n^{1/6}$.

For fixed $\beta>1$ at zero field, the Gaussian Poincar\'e inequality implies
that $\Var(F_n)\le C_\beta n$. Chatterjee improved this to $O_\beta(n/\log n)$ and
related this superconcentration to disorder
chaos~\cite[Theorem~1.5]{chatterjee-2009}. Aronow and Lopatto proved
$c_\beta n^{4/15}\le\Var(F_n)\le C_v n^{7/15}$ for all sufficiently large
$n$~\cite[Theorem~1.4]{quantitative-parisi}, using quantitative comparisons
between positive fractional moments of $Z_n$ and constrained Parisi
functionals~\cite[Theorem~1.1]{quantitative-parisi}. The sharp exponent was
then obtained by OpenAI~\cite[Theorem~1.1]{oai-fluctuations}, as stated in
\cref{f:main}.

\Cref{f:main} is a statement about exponents. It identifies neither a
limiting law nor a positive limit of $n^{-1/3}\Var(F_n)$. Both are stated in
the companion preprint~\cite[Theorems~1.1 and~1.2]{oai-limiting-law}: there,
$n^{-1/3}\Var(F_n)$ converges to a constant $a_2(\beta)\in(0,\infty)$, and
$(F_n-\E F_n)/\sqrt{\Var(F_n)}$ converges in law to a nondegenerate
probability measure with mean zero and variance one. The proofs
in~\cite{oai-limiting-law} use the exponent of \cref{f:main} and the
comparison estimates of~\cite{oai-fluctuations};
see~\cite[Appendix~A]{oai-limiting-law}. We do not use~\cite{oai-limiting-law} in
this paper. On the side of large deviations, Pei identified the coefficient
of the fifth-order term in the constrained Parisi functional and the
resulting rate for upward deviations of $F_n/n$, with $n\to\infty$ taken
before the deviation tends to zero~\cite[Theorem~1]{pei-fifth-order-2026}.

The proof of \cref{f:main} in~\cite{oai-fluctuations} reduces it to a
two-sided bound on a positive logarithmic moment. Let $F_n^c=F_n+\beta G$,
where $G$ is a centered Gaussian variable of variance $1/2$, independent of
the disorder. The variable $\beta G$ completes the covariance of $\beta H_n$ to
$n\beta^2R(x,y)^2/2$, and we call $F_n^c$ the completed log partition
function (\cref{sec:model-completion}). For every fixed $m\in(0,1/300)$ and
$s_n=n^{-1/6+5m}$, it is shown in~\cite[Proposition~7.5]{oai-fluctuations}
that $\log\E\exp\{s_n(F_n^c-\E F_n^c)\}$ lies between constant multiples of
$n^{20m}$ and $n^{100m}$ for all sufficiently large $n$. We state this bound
precisely as \cref{f:window}. The deduction of \cref{f:main} from this bound
in~\cite[Lemmas~7.6 and~7.7]{oai-fluctuations} uses an argument of Aronow
and Lopatto~\cite[Proposition~5.1]{quantitative-parisi}, based on the
log-concavity of an exponential augmentation of $F_n^c$. We reproduce it in
\cref{sec:fluct}.

\paragraph{The mean bias.}
For fixed $\beta>1$, Aronow and Lopatto proved that
$c_\beta n^{-11/12}\le\Delta_{n,\beta}\le C_\beta n^{-2/3}$ for all
sufficiently large $n$~\cite[Theorem~1.4]{quantitative-parisi}, and later
that $\Delta_{n,\beta}\le C_b n^{-7/10}$ for every
$n\ge2$~\cite[Theorem~1.1]{hexagon-bias}. Their lower bound $n^{-11/12}$
follows from \eqref{eq:intro-variance-deficit} and the lower bound
$c_\beta n^{4/15}$ on the variance; with \eqref{f:variance} in place of this
bound, the exponent becomes $5/6$, as in \cref{cor:mean-bias}. The proof of
\eqref{eq:intro-variance-deficit} balances a gain that is linear in a small
positive tilt and proportional to the variance against the fifth-order cost
of constraining the Parisi measure to have mass at least $s$ at
zero~\cite[Lemma~7.2]{quantitative-parisi}. In physics, the correction is
expected to be of order $n^{-a}$ with $a<1$ at low temperature, and
two values have been proposed. The value $a=2/3$ was supported by the
zero-temperature numerics of Bouchaud, Krzakala, and
Martin~\cite{bouchaud-krzakala-martin-2003} and of
Boettcher~\cite{boettcher-2005}, by Billoire's numerical integration along
the interpolation of Guerra and Toninelli~\cite{guerra-toninelli-2002}
in~\cite{billoire-2006}, and by the
argument of Aspelmeier, Billoire, Marinari, and Moore mentioned above.
Franz, Parisi, and Ricci-Tersenghi argued that the correct value is
$a=3/4$, possibly up to logarithmic
factors~\cite[equations~(13)--(17)]{franz-parisi-ricci-2026}. The bound
of~\cite{hexagon-bias} excludes a nonzero correction of order $n^{-2/3}$ at
every fixed $\beta>1$ and is consistent with the value $3/4$. A fuller account
of these predictions is given in~\cite[Section~1.4]{hexagon-bias}.

\paragraph{Normalizations.}
We have converted the results quoted in this section and the previous one to
the normalization \eqref{eq:intro-model}. Some of the cited works sum over
ordered pairs of sites, which changes the inverse temperature by a constant
factor; for instance, the critical inverse temperature
of~\cite{chen-lam-2019} is $1/\sqrt2$ in its own normalization.
\Cref{sec:model-dictionary} compares our normalization with those of the
works on which our proofs rely.

\section{Formal verification}\label{sec:intro-formal}

OpenAI has published formalizations of~\cite{oai-barriers}
and~\cite{oai-fluctuations} in the Lean~4 proof
assistant~\cite{lean4,mathlib,oai-math-lean}. The main formal theorem on the
fluctuation exponents takes the two-sided bound of \cref{f:window} as an
explicit hypothesis. In~\cite{oai-fluctuations}, that bound is the
conclusion of Proposition~7.5, which is proved there on paper. We have
completed this formal proof: our Lean development proves the three
conclusions of \cref{f:main} with no hypothesis beyond the definition of the
model and $\beta>1$. It follows the mathematical route
of~\cite{oai-fluctuations}; the work consisted in carrying out that route in
full detail for the comparison problems of~\cite[Section~3]{oai-fluctuations}
whose minimizers have a positive lower radius, and for the successive
comparisons built from them.

We have also formalized \cref{thm:barrier-730}, \cref{b:cor:subcritical},
and \cref{cor:mean-bias}. In these formal proofs, the upper bound on the
mean bias of~\cite[Theorem~1.1]{hexagon-bias} and the variance--bias
inequality \eqref{eq:intro-variance-deficit} are proved, not assumed;
their formal proofs come from our Lean development~\cite{hexagon-bias-lean}
for~\cite{hexagon-bias}. The formal proof of \cref{cor:mean-bias} also
proves the identities that relate the normalizations of \cref{f:main} and
\eqref{eq:intro-variance-deficit}. The formal proof of
\cref{thm:barrier-730} has the structure of the written proof, but it uses
a weaker form of the local expansion of \cref{b:prop:os-expansion}, which
suffices for the same exponent. \Cref{thm:barrier-strong} is not
formalized.

Each of these four formal theorems was compared with a separately written
statement of the same result that depends only on Mathlib, and its proof was
checked to use only the standard axioms \texttt{propext},
\texttt{Classical.choice}, and \texttt{Quot.sound}. The declarations on
which its statement and proof depend were then checked again by three
kernels: the kernel of Lean and the independent kernels nanoda and con-ron.
\Cref{sec:formal} states the formal theorems, explains how they correspond
to the statements of this paper, describes what the completion of the proof
of \cref{f:main} had to establish, and describes these checks.

\statusIntroFormal{} The proofs in this paper are complete without reference
to the formalization.

\section{Methods and their sources}\label{sec:intro-methods}

We describe the proofs of \cref{thm:barrier-730,thm:upper} and the sources
of their ideas. A more detailed overview is given in \cref{sec:overview}.

\subsection{The lower bound}

\paragraph{From overlap constraints to slow mixing.}
We follow the strategy of~\cite{oai-barriers}. Its first step reduces the
problem to a statement about the stationary
chain~\cite[Section~8.1]{oai-barriers}. Let $v$ be a reference
configuration with $R(v,X_0)$ bounded away from zero. By spin-flip symmetry,
the half-space $\{y:R(v,y)\le0\}$ has Gibbs mass at least $1/2$. If
$d^J(X_0,t)\le\tau<1/2$, the chain started at $X_0$ is in this
half-space at time $t$ with probability at least $1/2-\tau$. A single update
changes $R(v,\cdot)$ by at most $2/n$, so the chain must cross every
intermediate overlap level. It suffices to show that the chain started from
$X_0\sim\pi^J$ crosses these levels before time $t$ only with small
probability. At each deterministic time the stationary chain has law
$\pi^J$, and a union bound over times controls the discrete chain; an
independent Poisson clock of rate $n$ then gives the continuous-time chain.
This argument is related to bottleneck estimates for Markov chains
(see~\cite[Theorem~7.4]{levin-peres-2017}). Here, however, the bottleneck
depends on the initial configuration, and it must separate almost every
initial configuration in Gibbs measure from equilibrium.

To show that crossings are rare, we build, as in~\cite{oai-barriers}, a
growing family of reference configurations along the path. The references
are chosen from independent Gibbs samples, and their overlaps with the
current configuration are constrained simultaneously. At each crossing
level, we find a new reference with large overlap with the current configuration
using a bound on the Gibbs mass of overlap
neighborhoods~\cite[Proposition~6.1]{oai-barriers}. This bound is derived
from an upper tail bound for the energy and an estimate for small Gaussian
rotations of the disorder~\cite[Lemmas~5.1 and~5.2]{oai-barriers}. We
then search backward for the crossing of the overlap level of the new
reference, while overlap constraints keep the overlaps with the earlier
references almost unchanged. Iterating this step creates many simultaneous
constraints at a single time. We arrange the references in blocks of fixed
length, so that a conflict deletes only a short final segment of a block, and
we draw the overlap thresholds at random, one level at a time; these features
allow us to obtain the final
probability estimate~\cite[Sections~8.3--8.5]{oai-barriers}.

The overlap constraints have the following form. Let $\varepsilon$ be a
tolerance. If the mean $r$ of the overlaps $R(v,x)$ and $R(v,w)$ lies in a
fixed compact subinterval of $(0,q_\beta)$ and $R(w,x)\ge r+\rho$, then
$\abs{R(v,x)-R(v,w)}$ does not lie in $[2\varepsilon,4\varepsilon]$. In words, a large overlap between $w$ and $x$
forces their overlaps with $v$ to agree. Since overlaps change by at most
$2/n$ in one update, excluding this band keeps the two overlaps close along
a path. To show that such a constraint fails with small Gibbs probability, we show
that the partition function restricted to the triples $(v,x,w)$ with
violating overlaps is small compared with $Z_n^3$. In physics, free energies constrained by the overlap
with an equilibrium reference configuration go back to Franz and
Parisi~\cite{franz-parisi-1995}. The work of Ben Arous and
Jagannath~\cite{ben-arous-jagannath-2018} is an earlier rigorous argument
that passes from such constrained free energies to the dynamics, there for
the spectral gap.

\paragraph{Constrained comparisons and the scalar problem.}
We bound constrained partition functions by Gaussian interpolation.
Guerra's hierarchical interpolation~\cite{guerra-2003} gives the upper bound
in the Parisi formula. Talagrand extended it to systems of replicas with
prescribed overlaps, in the proof of the matching lower
bound~\cite[Section~3]{parisi-formula} and in the two- and three-replica
bounds of~\cite[Propositions~3.4 and~4.4]{talagrand-2007}. These bounds
already retain a cost when the prescribed overlaps differ from the endpoint
of the interpolation. Bovier and Klimovsky developed the interpolation for
multidimensional spins~\cite{bovier-klimovsky-2009}. Chen studied the
two-dimensional Guerra--Talagrand bound and its variational
representation~\cite{chen-2017}, and Panchenko proved the vector-spin
variational formula and the asymptotic sharpness of the constrained
bounds~\cite{panchenko-2018}. We use the matrix comparison for finite hierarchies
of~\cite[Proposition~2.3]{oai-barriers} for the quadratic covariance of SK,
which allows any finite number of replicas and endpoint matrices whose
diagonal entries need not equal one. The hierarchies themselves descend from Parisi's
recursion~\cite{parisi-1979} and from Ruelle's probability
cascades~\cite{ruelle-1987}; see~\cite{panchenko-2013}. For the mixed even
$p$-spin models with exponentially summable coefficients, including SK,
Aronow and Lopatto identified the limiting overlap array averaged over the
disorder~\cite[Theorems~1.1 and~1.2]{al-overlap-structure}. At zero field, the
absolute overlaps have the law of the overlap array $(U_{ab})_{a<b}$ of the
corresponding Ruelle probability cascade, and the signed array has the law of
$(\epsilon_a\epsilon_bU_{ab})_{a<b}$, where the replica signs $\epsilon_a$
are independent symmetric signs, independent of the cascade. These are
limiting statements, while the proof below needs estimates at finite $n$ on
shrinking scales.

We obtain the pressure deficits from the scalar Parisi problem. Auffinger and
Chen proved that the Parisi functional has a unique minimizer, and they
expressed it through a stochastic control
problem~\cite[Theorems~2 and~3 and Corollary~1]{auffinger-chen-uniqueness-2015}.
They also proved regularity properties of the Parisi equation and of its
minimizers~\cite[Propositions~1 and~2 and Theorems~1 and~2]{auffinger-chen-2015}.
Jagannath and Tobasco gave a dynamic-programming treatment of the
functional~\cite[Lemma~18]{jagannath-tobasco-2016} and characterized its
minimizers through their first
variation~\cite[Proposition~1.1 and Corollary~3.10]{jagannath-tobasco-2017};
Chen obtained the same characterization~\cite[Theorem~2 and
Proposition~1]{chen-2017}. Toninelli proved that replica symmetry is broken
at zero field for $\beta>1$~\cite{toninelli-2002}. Let $u$ be the spatial
derivative of the solution of the Parisi equation at the minimizer, and let
$Y$ be the associated diffusion. The identity $\E u(t,Y_t)^2=t$ at every
point $t$ of the support is due to
Talagrand~\cite{talagrand-2006-parisi-measures}; see
also~\cite[Theorem~5]{auffinger-chen-2015}, \cite[Proposition~1]{chen-2017},
and~\cite[Lemma~2.2]{full-rsb}. OpenAI proved it in the normalization used
here, together with the regularity of the scalar evolution that it
requires~\cite[Lemma~3.1 and Proposition~3.2]{oai-barriers}; the argument
of~\cite{oai-barriers} does not assume a description of the support, and it
treats separately the case of a forbidden overlap
interval~\cite[Section~8.2]{oai-barriers}. By~\cite[Theorem~1.1]{full-rsb}, the identity holds on the whole
interval $(0,q_\beta)$, and this case does not arise. Differentiating the
identity once and twice along this interval, we obtain two stationarity
identities,
\[
 \beta^2\E\chi(t,Y_t)^2=1,\qquad
 \E w(t,Y_t)^2=2\alpha(t)\E\chi(t,Y_t)^3\qquad(0<t<q_\beta),
\]
where $\chi$ and $w$ are the second and third spatial derivatives and
$\alpha$ is the distribution function of the Parisi measure
(\cref{sec:scalar}). Both identities appear, in the normalization of energy
units, in the proof of~\cite[Lemma~7.2]{quantitative-parisi}
(see~\cite[equation~(7.12)]{quantitative-parisi} and the identity preceding
it), and the differentiated moment identity is also
in~\cite[Lemma~2.5]{full-rsb}. The first identity expresses, in scalar form, the
marginal stability of the Parisi solution, which De Dominicis and Kondor
found in a truncated replica model~\cite{de-dominicis-kondor-1983}; Crisanti
and Rizzo wrote it in physical
variables~\cite[equations~(22) and~(23)]{crisanti-rizzo-2002}.

Positive fractional moments $\E Z_n^s$ have their own genealogy. Talagrand
computed the limits of their normalized logarithms~\cite{talagrand-2007-large-deviations}.
For SK at zero field and fixed $s\in(0,1)$, the limit of
$(ns)^{-1}\log\E Z_n^s$ is the minimum of the Parisi functional over measures
with mass at least $s$ at zero. This formula is contained in the multiscale
formula of Contucci and Mingione~\cite[Theorem~2.1]{contucci-mingione-2019},
and it was stated in this form by Chen, Guionnet, Ko, Lacroix-A-Chez-Toine,
and Mourrat~\cite[Theorem~2.1]{chen-guionnet-ko-lacroix-mourrat-2026}. Aronow
and Lopatto made the comparison quantitative when $s$ depends on
$n$~\cite[Theorem~1.1]{quantitative-parisi}, and they proved that the
constraint raises the Parisi value by an amount of order
$s^5$~\cite[Lemma~7.2 and Proposition~7.3]{quantitative-parisi}. This is a
rigorous counterpart of the replica calculations of Kondor and of Parisi and
Rizzo, and Pei identified its leading
coefficient~\cite[Theorem~1(a)]{pei-fifth-order-2026}.

\paragraph{What is new.}
Four parts of the proof are new. Together they allow overlap constraints at
the scale $\rho\asymp n^{-7/30}$, with tolerance comparable to $\rho$.
\begin{enumerate}
\item \emph{A three-replica covariance perturbation with a cubic deficit}
(\cref{sec:cubic}). In~\cite[Section~7]{oai-barriers}, the overlap
constraints near a level $r$ come from a local perturbation of the covariance
path of width $\rho^8$ with tolerance $\rho^{20}$, and the resulting pressure
deficit is of order $\rho^{43}$. We keep the block hierarchy and the
antisymmetric Gaussian contrast of~\cite[Section~7, Steps~1 and~2]{oai-barriers},
but we perturb the covariance path differently. It is positive
semidefinite at every stage, ends exactly at the prescribed overlap matrix
after a terminal correction, and its expansion uses both stationarity
identities to cancel the lower-order terms. The resulting
deficit is of order $\rho^3$ per spin, with tolerance comparable to $\rho$,
uniformly for $r$ in a compact subinterval of $(0,q_\beta)$. Its coefficient
involves $\E\chi(r,Y_r)^3$, which is positive (\cref{sec:scalar}).
\item \emph{A positive fractional-moment comparison} (\cref{sec:fractional}).
For constraints near the level $r$, the deficit $n\rho^3$ in the logarithm
of a constrained partition function is smaller than the scale $O_\beta(\sqrt n)$
supplied by the standard Gaussian concentration bound for Lipschitz
functions, so that bound is too coarse here. Instead, we bound the moment $\E Z^s$ of a constrained
partition function $Z$, with $s$ comparable to $\rho$, by a fractional-moment
version of the matrix comparison of~\cite[Proposition~2.3]{oai-barriers},
in which the outer level of the hierarchy has mass $s$, as in the
positive-moment interpolation of~\cite[Section~2]{quantitative-parisi}. The
comparison uses the profile $\max\{\alpha,3s\}$, whose cost is of order
$s^5$ by the trial bound of~\cite[Lemma~7.2]{quantitative-parisi}, together
with a bound on the response of the local perturbation to this change of
profile. Markov's inequality then bounds the disorder probability of
failure of each of these bounds on constrained partition functions by
$e^{-cn\rho^4}$.
\item \emph{One disorder event for all constraints} (\cref{sec:static}). The
denominator $Z_n^3$ is bounded below on the event
$F_n/n\ge P_\beta-c_0\rho^3$, for a fixed $c_0>0$. Let $\delta\in(0,c_0)$ be
fixed such that $\Delta_{n,\beta}\le(c_0-\delta)\rho^3$ for all
sufficiently large $n$; this margin is obtained by choosing the constant $R$
in $\rho=Rn^{-7/30}$ large. Chebyshev's inequality and the variance bound
of~\cite[Theorem~1.4]{quantitative-parisi} give
\[
 \PP\bigl\{F_n/n<P_\beta-c_0\rho^3\bigr\}
 \le\frac{\Var(F_n)}{\delta^2n^2\rho^6}
 \le C_{\beta,\delta}\,n^{-23/15}\rho^{-6},
\]
which is $O(n^{-2/15})$ for $\rho=Rn^{-7/30}$. We intersect this event with
the bounds on polynomially many constrained partition functions, and the
complement of the resulting event, whose probability tends to zero, is
discarded once and for all. On this event, the number
of reference configurations and path steps, which is exponential in
$1/\rho$, multiplies only the conditional Gibbs probability of a violation,
which is at most $e^{-cn\rho^3}$. This replaces the probability estimate
of~\cite[Lemma~4.2]{oai-barriers}.
\item \emph{Control of crossing errors through the next retained reference}
(\cref{sec:induction,sec:barrier-proof}). In~\cite[Section~8.4]{oai-barriers},
the overlap of an earlier reference with the current configuration may drift
by up to $6j\varepsilon$ after $j$ levels, and the final count requires
$p\varepsilon/\rho$ to be small for $p$ levels. Since $p$ is of order
$1/\rho$, this forces $\varepsilon$ to be much smaller than $\rho^2$.
Instead, we control the overlap of each retained reference with the next
retained reference. This overlap changes only when a deletion occurs, and the
error stays of order $\varepsilon$ at every level, rather than accumulating.
The first reference of each block keeps its next retained
reference, and at one common time the overlap of the path with the first
reference of each of the $M$ blocks lies within a fixed multiple of
$\varepsilon$ of its threshold. A count over the $m+1$ path steps of these
$M$ windows of predictable center gives the bound $(m+1)\gamma_n^M$, where
$\gamma_n$ is the ratio of the window length to the length of a threshold
interval, and $\gamma_n$ stays below one when $\varepsilon$ is a small fixed
multiple of $\rho$. This allows $\varepsilon$ comparable to $\rho$.
\end{enumerate}
The other overlap constraints use the inherited pressure bounds:
the pair bound of~\cite[Lemma~4.3]{oai-barriers}, the three-replica bound at
fixed scales of~\cite[equation~(4.20)]{oai-barriers}, and, for large
overlaps $R(w,x)$, the comparison of~\cite[Section~7, Steps~5--8]{oai-barriers}.
We also extend the bound on the Gibbs mass of overlap neighborhoods
of~\cite[Proposition~6.1]{oai-barriers} to every scale $L\to\infty$ with
$L=o(n)$; its proof is that of~\cite{oai-barriers}
(\cref{sec:references}).

\paragraph{The exponent.}
We determine the exponent by balancing these estimates. The lower
bound on $Z_n^3$ in item~3 requires the mean bias $\Delta_{n,\beta}$ to be
smaller than a small multiple of the deficit $\rho^3$. The comparison proved
in~\cite[Theorem~2.4]{oai-barriers} has error $O(n^{-1/24})$. This is much
smaller than the deficit $\rho^{43}$ at the scale $\rho=n^{-1/10000}$ used
there, but not than $\rho^3$ at the scale used here. For this reason we use
the upper bound $\Delta_{n,\beta}\le C_b n^{-7/10}$ of Aronow and
Lopatto~\cite{hexagon-bias}, and we take $\rho=Rn^{-7/30}$ for a large
constant $R$. A fixed overlap interval inside $(0,q_\beta)$ then contains
of order $1/\rho\asymp n^{7/30}$ crossing levels. By counting, we find
that the probability of crossing all of them before time $t$ is small
when $\log t$ is a small multiple of the number of levels. This leads to the
time $\exp(A_\beta n^{7/30})$. Two further conditions concern the two kinds
of union bound in the proof. For the union over reference configurations and
path steps, the conditional Gibbs error is at most $\exp(C/\rho-cn\rho^3)$,
which tends to zero if $cn\rho^3-C/\rho\to\infty$. For the disorder event,
a union over at most $Cn^d$ constrained partition functions has failure
probability at most $Cn^de^{-c'n\rho^4}$, which tends to zero if
$c'n\rho^4-d\log n\to\infty$. The
condition $n\rho^4/\log n\to\infty$ implies both. For $\rho=Rn^{-\lambda}$
with fixed $R$, it holds when $\lambda<1/4$, in particular at
$\lambda=7/30$; at $\lambda=1/4$, the second union bound fails. This is the
restriction to exponents below $1/4$ mentioned in \cref{sec:intro-results}.

\paragraph{Inputs for the lower bound.}
Besides standard tools, the proof of \cref{thm:barrier-730} uses the following
public results, which are stated in \cref{sec:inputs} with their hypotheses:
the Parisi formula and Guerra's bound~\cite{guerra-2003,parisi-formula}, the
uniqueness of the Parisi minimizer~\cite{auffinger-chen-uniqueness-2015}, the
interval support of the Parisi measure~\cite{full-rsb}, the upper bound on
the mean bias of~\cite[Theorem~1.1]{hexagon-bias}, the upper bound on the
variance and the trial bound of~\cite{quantitative-parisi}, and results
of~\cite[Sections~2--8]{oai-barriers}. The proof of
\cref{thm:barrier-strong} uses in addition the small-deviation inequality
of~\cite{paouris-valettas-2018}. All other steps are proved in this paper.

\subsection{The upper bound}

\paragraph{From flows to mixing.}
We follow the strategy of~\cite{oai-typical-upper}. Its last step reduces the
problem to a lower bound on stationary flows across
cuts~\cite[Section~7, Lemma~8]{oai-typical-upper}. For $S\subseteq\Sigma_n$,
the flow $\sum_{x\in S,\,y\notin S}\pi^J(x)\mathcal L^J(x,y)$ is the stationary
rate at which the continuous heat-bath process jumps from $S$ to its
complement. Let $0<\epsilon<1/3$ and $\phi>0$, and suppose that the flow
across $S$ is greater than $\phi$ for every $S$ with
$\epsilon\le\pi^J(S)\le1-\epsilon$. Then for every $t>0$,
\[
 \sum_{x\in\Sigma_n}\pi^J(x)\,d^J(x,t)\le\frac{n\sqrt{2\log2}}{\phi\sqrt t}+2\epsilon.
\]
The proof removes a set of small Gibbs mass, proves a Poincar\'e inequality
on the remainder by the coarea argument for the reversible Cheeger
bound~\cite[Theorem~3.5]{lawler-sokal-1988}, and then bounds the Gibbs
average of the distances through the entropy dissipation of the whole
chain~\cite[Lemmas~2.5 and~2.7]{diaconis-saloff-coste-1996}. It uses only
cuts that separate appreciable Gibbs masses, and it does not require the
trajectory to stay in the retained set. We reproduce it in
\cref{u:lem:pruning}. With $\epsilon=n^{-\mathsf d}$,
$\phi=e^{-n^{1-\mathsf q}}$, and $t=e^{n^{1-\omega}}$, where $\mathsf d$ and
$\mathsf q$ are positive exponents of the proof (\cref{u:sec:parameters}),
the right side tends to zero when $\omega<\mathsf q$, and using Markov's
inequality for the Gibbs average, we obtain \cref{thm:upper}. It therefore suffices to show
that, with probability tending to one over the disorder, the flow across
every $S$ with $\min\{\pi^J(S),\pi^J(S^{\rm c})\}\ge n^{-\mathsf d}$ is greater
than $e^{-n^{1-\mathsf q}}$.

\paragraph{The proposal.}
To bound these flows, OpenAI constructs a random configuration correlated
with the disorder, called a proposal~\cite[Section~4]{oai-typical-upper}.
Each coordinate of the proposal follows the magnetization martingale of the
scalar Parisi diffusion. The martingale is driven by successive queries of
the interaction matrix: over each cell of a time mesh, the increment of the
magnetizations, orthogonalized against the earlier query directions, is
multiplied by the matrix. By orthogonal invariance, the part of the matrix
that has not been revealed is a compressed block of the Gaussian orthogonal
ensemble (GOE) at each step. Beyond the top $q_\beta$ of the support of the
Parisi measure, the martingale is continued until the average uncertainty of
the coordinates is small, and the query noise is mixed with an independent
Brownian noise. At the final time, the spins are drawn independently with
the magnetizations as means. Successive Gaussian conditioning along
orthogonalized directions is used in Bolthausen's iterative construction of
solutions of the TAP equations~\cite[Sections~2--3]{bolthausen-2014}, and
Montanari's incremental approximate message passing algorithm uses a Parisi
diffusion and normalized increments of the
magnetizations~\cite[Sections~2--3]{montanari-2021}.

The equilibrium input is the stationarity of the Parisi measure along the
whole interval $[0,q_\beta]$. The stationarity conditions on the support were
obtained by Chen~\cite[Theorem~2 and Proposition~1]{chen-2017} and by
Jagannath and Tobasco~\cite[Proposition~1.1 and
Corollary~3.10]{jagannath-tobasco-2017}, and the support is the whole interval
by~\cite[Theorem~1]{zhou-2026} for $\beta$ close to one and
by~\cite[Theorem~1.1]{full-rsb} for every $\beta>1$. Along the scalar
diffusion, an exact identity balances the information about the matrix that
the queries reveal against the energy gained by the
proposal~\cite[Section~3, Lemma~4]{oai-typical-upper}. It yields the two
estimates of~\cite[Proposition~5]{oai-typical-upper}. The relative entropy of
the joint law of the matrix and the proposal with respect to the joint law of
the disorder and a Gibbs sample is sublinear in $n$, and outside an event of
small probability, the law of the pair is dominated by $e^{o(n)}$ times this
joint law.

\paragraph{Coverage and the connecting sequence.}
Forward relative entropy alone does not show that the proposal reaches a
prescribed set of positive Gibbs mass. Under the joint law of the disorder
and a Gibbs sample, the conditional law of the disorder given the spin is
strongly log-concave. A transport bound for such
laws~\cite[Corollary~3.10]{bogachev-kolesnikov-medvedev-2005} and Gaussian
smoothing of the disorder~\cite[Lemma~11]{lugosi-neu-2022} reverse the entropy
comparison. They give a lower bound on the probability that the proposal lands
in any set of appreciable Gibbs mass~\cite[Section~5,
Proposition~6]{oai-typical-upper}. Two proposals generated from a common
smoothed disorder then land on the two sides of a balanced cut with a
probability that is not too small. Following~\cite{oai-typical-upper}, these
two proposals are run to a longer, fine horizon, and they are joined by a
sequence of proposals run to a shorter, coarse horizon, obtained by rotating
the Gaussian inputs, with small consecutive Hamming distances. Near-pinning of
the magnetizations joins the two horizons. If the cut had small flow, its
enlarged vertex boundary would have tiny Gibbs mass, and the likelihood bound
keeps the coarse proposals away from this boundary. But a sequence with short
steps that crosses the cut must come close to its boundary. This
contradiction bounds the flows from below~\cite[Section~6,
Proposition~7]{oai-typical-upper}.

\paragraph{What is new.}
In~\cite{oai-typical-upper}, the parameters of the proof are fixed numbers,
and the saving $1/40000000$ is not optimized~\cite[Section~1, after
Theorem~1]{oai-typical-upper}. To reach every saving below $1/4$, the
estimates for the coarse proposal must hold outside events of probability
close to $e^{-n^{3/4}}$, with a likelihood bound close to $e^{n^{3/4}}$ times
the joint law. The following refinements yield every fixed saving below $1/4$.
\begin{enumerate}
\item \emph{Sharp estimates for the coarse proposal}
(\cref{u:sec:concentration,u:sec:coarse}). We strengthen the scalar
entropy--energy identity of~\cite[Section~3, Lemma~4]{oai-typical-upper}
to a pathwise terminal identity. Along the scalar diffusion beyond the
support, the information revealed by the queries and $\beta$ times the
linearized energy of the proposal cancel up to a terminal term. The centered fluctuation
of this term has variance at most $\beta^4p(V)^3$, where $p(V)$ is the mean
conditional variance of a spin given its magnetization at the final time $V$
(\cref{u:lem:joint}). The coarse likelihood estimate absorbs the
deterministic amplitude of this term. The centered difference between this cancellation
and its version on the mesh has a moment generating function bounded by
$e^{C_\beta\theta^2h}$ for $\abs\theta\le c_\beta/h$, where $h$ is the mesh
parameter (\cref{u:lem:lag}). We combine these bounds with concentration of
weighted dyadic sums over cells of unequal length, normalized by the
durations of the cells (\cref{u:lem:weighted}), with a bound on the deviation
of the Gram matrix of the queries from the identity that uses the variance of
the squared projections and not only their tails (\cref{u:lem:gram}), and
with sharper bounds on the variances of the Gaussian terms that the proposal
does not use (\cref{u:lem:weighted-energy}). The dyadic triangular and
Cholesky arguments build on~\cite[Section~4, proof of
Proposition~5]{oai-typical-upper}.
\item \emph{The growing mesh} (\cref{u:sec:proposal}). The cells of the mesh
in~\cite{oai-typical-upper} have a common length. Ours have length $he^{t/4}$
at time $t$, with $h=n^{-\mathsf a}$, so that the horizon $L\log n$ is
reached after at most $8n^{\mathsf a}$ cells when $\mathsf a>L/4$
(\cref{u:lem:quad}). The quadrature error is $O_\beta(h)$, uniformly in the
horizon. On this mesh, adjacent cells have different lengths, and
the estimates distinguish the duration of the previous cell, which governs
the conditional law of the spins, from that of the next cell, which
normalizes the query. For instance, the bound on the conditional
correlations of the query entries with the spins in \cref{u:sec:fine-energy}
loses a factor two compared with equal cells.
\item \emph{Direct expected energy at the fine horizon}
(\cref{u:sec:fine-energy}). We compute the expected queried energy through
the exact one-query conditional identity of \cref{u:lem:onequery}, including
fallback outcomes. This avoids the Gram-event restriction in the energy
calculation of~\cite[Section~4, Proposition~5]{oai-typical-upper} and gives
the fine relative-entropy bound of \cref{u:prop:fine-kl}, with saving
$\mathsf k_{\rm f}<\min(\mathsf a,\mathsf f,1-\mathsf a)$.
\item \emph{The parameter system} (\cref{u:sec:budget}). The exponents of the
proof are affine functions of $\varepsilon_\omega=(1/4-\omega)/20$, and this
single family satisfies every condition of the proof for every
$\omega\in(0,1/4)$ (\cref{u:prop:range}).
\end{enumerate}

The connecting sequence, including fine proposals only at the two
endpoints, stability from coarse Gram events, and pinning without a fine
Gram exception, follows~\cite[Section~6, proof of
Proposition~7]{oai-typical-upper}.

\paragraph{The exponent.}
The proof of \cref{u:prop:range} shows that, for every choice of the
exponents, the estimates of \cref{part:upper} require
\[
 \omega<\mathsf q<\mathsf s_{\rm b}<\mathsf s_{\rm seq}<\mathsf e<\mathsf u_{\rm E}
 <\mathsf k_{\rm f}/2\le1/4 .
\]
The mixing step needs $\omega<\mathsf q$, since the flow threshold is
$e^{-n^{1-\mathsf q}}$. In the connecting sequence, $e^{-n^{1-\mathsf s_{\rm b}}}$
bounds the failure probability of one coarse proposal,
$e^{-n^{1-\mathsf s_{\rm seq}}}$ that of the whole sequence, and
$e^{-2n^{1-\mathsf e}}$ bounds from below the probability that the two fine
endpoints lie on opposite sides of the cut, and the proof needs
$\mathsf q<\mathsf s_{\rm b}<\mathsf s_{\rm seq}<\mathsf e$. Smoothing and
coverage need $\mathsf e+\mathsf d<\mathsf u_{\rm E}
<\min\{\mathsf D,\mathsf k_{\rm f}-\mathsf D\}$, where $n^{-\mathsf D}$ is the
smoothing parameter and $n^{1-\mathsf k_{\rm f}}$ bounds the forward relative
entropy at the fine horizon. Finally, this relative entropy contains terms of
order $nh=n^{1-\mathsf a}$ and $n^{\mathsf a}$, the second from the number of
queries, so that $\mathsf k_{\rm f}<\min\{\mathsf a,1-\mathsf a\}\le1/2$.

\paragraph{Inputs for the upper bound.}
Besides standard tools of Gaussian analysis and stochastic
calculus~\cite{tsirelson-ibragimov-sudakov-1976,karatzas-shreve-1991,boucheron-lugosi-massart-2013},
the proof of \cref{thm:upper} uses the following results, which are stated in
\cref{sec:inputs} and \cref{u:sec:scalar} with their hypotheses: the Parisi
formula and Guerra's bound (\cref{inp:parisi}), the uniqueness of the Parisi
minimizer (\cref{inp:unique}), the interval support of the Parisi measure
(\cref{inp:full-rsb}) and the identity $\Gamma(t)=t$ on $(0,q_\beta)$
(\cref{cor:inp-identity}), the regularity and scalar evolution of
\cref{inp:ob-scalar}, the scalar estimates of~\cite[Section~3, Lemmas~3
and~4]{oai-typical-upper}, and the transport, smoothing, and finite-chain
results cited above. All other steps are proved in this paper.

\subsection{The fluctuation theorem}

The proof of \cref{f:window} in~\cite{oai-fluctuations} combines finite
replica trees and covariance differentiation for Gaussian
hierarchies~\cite[Section~2]{oai-fluctuations}, optimized short comparisons
over overlap quantiles~\cite[Section~3]{oai-fluctuations}, a one-site cavity
estimate, which compares an overlap moment of a comparison with the
corresponding moment of a scalar recursion~\cite[Proposition~4.1]{oai-fluctuations},
a rigidity theorem, by which rescaled overlap quantiles of minimizing
comparisons have linear limits at small masses~\cite[Theorem~5.1]{oai-fluctuations},
finite propagation of regularity through successive
comparisons~\cite[Proposition~6.6]{oai-fluctuations}, and a low-field
perturbation that yields the upper bound~\cite[Section~7]{oai-fluctuations}.
As explained in~\cite[Section~1]{oai-fluctuations}, its closest analytic
predecessor is the combination in~\cite{quantitative-parisi} of a pinned
Brascamp--Lieb estimate for the variance~\cite{brascamp-lieb-1976}, optimized
covariance interpolation, and positive
moments~\cite[Propositions~2.2 and~2.3 and Section~3]{quantitative-parisi}.
The single-site analysis of Thouless, Anderson, and
Palmer~\cite{thouless-anderson-palmer-1977} is an early antecedent of the cavity
method in spin glasses, and M\'ezard, Parisi, and Virasoro developed
the cavity construction of the solution with broken replica
symmetry~\cite{mezard-parisi-virasoro-1986}. The cavity variational principle
of Aizenman, Sims, and Starr~\cite{aizenman-sims-starr-2003} is a rigorous
antecedent of comparisons that add or remove sites. \Cref{sec:fluct} summarizes this route, with
references to the corresponding statements.

\section{Organization of the paper}\label{sec:intro-organization}

\Cref{sec:overview} gives an overview of the proofs of both bounds.
\Cref{part:prelim} contains preliminaries: \cref{sec:model} defines the model,
the dynamics, and the joint law of the disorder and the initial
configuration, and \cref{sec:inputs} defines Gaussian hierarchies and the
Parisi functional and states the imported results. \Cref{part:static}
contains the static estimates for the proof of \cref{thm:barrier-730}.
\Cref{sec:scalar} treats the scalar
problem near a point of the support of the Parisi measure,
\cref{sec:flow} proves short-time estimates for Gaussian hierarchies,
\cref{sec:cubic} constructs the three-replica covariance perturbation and
proves the cubic deficit, \cref{sec:fractional} proves the fractional-moment
comparison, and \cref{sec:static} combines these into a single estimate for
all overlap constraints. \Cref{part:dynamics} contains the dynamical part of
the proof of the lower bound. \Cref{sec:references} treats the Gibbs mass of
overlap neighborhoods and the families of reference configurations,
\cref{sec:induction} proves that the overlap constraints persist along a
trajectory, and \cref{sec:barrier-proof} completes the proof of
\cref{thm:barrier-730} and \cref{b:cor:subcritical} and proves
\cref{thm:barrier-strong}.

\Cref{part:upper} proves the upper bound. \Cref{u:sec:setup} introduces the
Gaussian matrix from which the disorder is built, the scalar Parisi process
and its continuation beyond the support, and the exponents of the proof.
\Cref{u:sec:proposal} constructs the Gaussian-query proposal on a growing
mesh, and \cref{u:sec:fine} bounds its forward relative entropy at the fine
horizon. \Cref{u:sec:concentration} proves the concentration estimates on the
growing mesh, and \cref{u:sec:coarse} proves the joint cancellation and the
likelihood bound for the coarse proposal. \Cref{u:sec:connections} proves the
coverage estimate and constructs the connecting sequence, and
\cref{u:sec:proof} proves \cref{thm:upper,cor:upper-discrete,cor:two-sided}
and verifies the conditions on the exponents.

\Cref{part:fluct} concerns the fluctuations and the mean bias:
\cref{sec:fluct} states the bound of \cref{f:window} and deduces
\cref{f:main} from it, and \cref{sec:bias} proves \cref{cor:mean-bias}.
\Cref{part:formal} contains \cref{sec:formal}, which describes the formal
proofs, and \cref{sec:open}, which lists open problems.
